\chapter{By disrupting that order-a way of surprising (I'm Not Afraid of Anything Anymore)}

If nothing else (and it is much else), Puella Magi Madoka Magica is a meticulously structured series. Every moment of it is carefully placed to advance a complex story with extensive character development in a surprisingly small space. One example of that meticulous structure is the symmetry throughout the series, the way in which moments large and small repeat at precisely the right time to reflect one another.

Two episodes utterly transform the series. On a first viewing, these two episodes stand out as moments when what appeared to be a show about one thing suddenly becomes a different show entirely. As in a mirror, they are at precisely opposite ends of the series, the third episode and the third-to-last. This is the first of these episodes.

The episode opens with hints to the second and longest of the series' three major arcs, focusing primarily on Sayaka and Kyoko. Sayaka is visiting a boy, Kyousuke Kamijo, to give him the CD she and Madoka went shopping for in the first episode. As the two listen to it together, we see Sayaka remember her first encounter with true beauty and with the human power of creation, when she saw Kyousuke perform when both were small. We never learn precisely what is wrong with Kyousuke, only that he has lost some mobility in his hands and can no longer play the violin; regardless, what matters is that the beauty Kyousuke could once create is now beyond his reach, trapped in the past. This is the earliest emergence in the series of one of its major themes, a clear marker of its Buddhist influences: the inevitability of decay and loss. The beauty Kyousuke creates had to be lost sooner or later, because everything is; time is the destroyer of all \cite{ref33}. The loss of Kyousuke's music is not qualitatively different from Mami's family's death in a car crash or the destruction of the universe by the unrelenting march of entropy; they differ only in scale.

But thankfully, the opening credits are here to save us from such melancholy thoughts! This is still the false Madoka Magica, after all, the safe, comfortable magical girl show with only occasional hints of darkness. Either the credits or Mami herself will always step in to save us before things get too dark.

But like teeth on the edges of the frame, darkness is creeping in around the show. Mami's flashback to her wish, to save her own life, contains so much unstated: Mami lives alone, with no clear source of money, and she was clearly in the backseat of the car that crashed. She wished hastily to live, and now she counsels Sayaka and Madoka to think carefully about their wishes and be absolutely certain they are wishing for what they want.

The strong implication is that she wished to live, when she could have wished for her whole family to live. The series' paratext suggests that much the same is true for Charlotte (or, at least, for the "witch of cheesecake" that appears to have been an early draft of Charlotte); she wished to share one last cake with her mother, when she could have wished for her mother not to die \cite{ref34}. 

But Charlotte is far from Mami's only parallel here. Last episode, we saw Mami's magical girl transformation paralleled with Junko's transformation into a different kind of warrior, the ambitious corporate climber. This episode, we see Junko laid low by an inevitable part of the life of the typical Japanese salaryman: the after work drunken bender \cite{ref35}. As she staggers into the Kaname home, she both foreshadows that Mami will shortly fall to an inevitable part of the life of a magical girl (namely, the messy death) and provides the impetus for a crucial conversation between Madoka and her father.

Madoka asks a natural question: why does her mother enjoy her life? What dream is she living out by being an ambitious cog in a profit machine? Her father explains that Junko's dream is not to do something, but to be something; that she works for the sake of working, that what she values about the effort is the effort itself. 

This mirrors the critical question Mami asks Sayaka. Does Sayaka wish to help Kyousuke, or to have helped Kyousuke? Does she want something for herself, in which case she should wish for that, or is it truly the helping itself that she wants? Just as Madoka is interested in being a magical girl, while Sayaka wants to fight evil, here Sayaka is focused on what she wants to do, rather than on what state of being she wants to achieve. She assumes that her action will bring that state of being about, but she is still failing to express the wish she truly desires. 

Similarly, Mami reveals in her final conversation with Madoka that she hates the state of loneliness in which she finds herself as a consequence of her wish. Though she stated earlier that she prefers how things are now to the prospect of death, she still regrets that she couldn't have made a better wish, and she still feels terribly, utterly alone. But just as she will at the end of the series, Madoka reminds Mami that she is not alone, and promises to become a magical girl to help support her. 

This is the moment at which Madoka kills Mami. The joy that Mami feels at knowing she is no longer alone causes her to showboat even more than last episode. She underestimates the threat Charlotte represents, and in so doing ensures her death. More importantly, just as with Kyousuke's music, her joy cannot last. It must end, decay, turn sour, because that is the inevitable consequence of existing within time.

Except for one thing: cheese.

We know from the paratext and from the Rebellion film that Charlotte is obsessed with cheese, searching for it endlessly. And what is cheese if not something good and life-sustaining that comes out of decay? It is rotten milk, raised into both a culinary delight and source of sustenance. It is a perfect example of the alchemical concept of putrefaction, the physical and spiritual notion that death is a source of life \cite{ref36}. Decay is repulsive, and yet the ugly, squirming mass of mold and maggots is teeming with life, able to sustain more beautiful and lovable creatures; without that decay, there would be no life. 

In devouring Mami, Charlotte finds her cheese. This death and decay brings forth a new life, because it is the moment at which Madoka Magica transcends the norms of its genre and begins to fulfill its potential. Only a few short minutes after Mami first attacks Charlotte, everything has changed: Mami is dead. Homura has saved Sayaka and Madoka. Kyubey offers no comfort as they sob in the hospital parking lot. And as "Magia" finally takes its place as the true ending credits, one thing is clear: Madoka Magica has begun.
